%==============================================================================
% CAPÍTULO 13 — TDD: TESTE ANTES DO MODELO
% PARTE III — Engenharia de Software Odontológico com SDD e TDD
%==============================================================================
\chapter{TDD: Teste Antes do Modelo}
\label{cap:tdd-python}
\niveltres

Se o SDD (Capítulo~12) estabelece \textbf{o quê} deve ser construído,
o \textit{Test-Driven Development} (TDD) define \textbf{como} verificar
que foi construído corretamente. Este capítulo apresenta o ciclo
Red--Green--Refactor aplicado ao desenvolvimento de software odontológico,
com ferramentas práticas do ecossistema Python e do pacote \texttt{odontoia}.

\section{Filosofia RED $\rightarrow$ GREEN $\rightarrow$ REFACTOR}

O TDD foi popularizado por Kent Beck no final dos anos 1990 e consolidado
em seu livro \textit{Test-Driven Development: By Example} (2003). A
filosofia é radicalmente simples — e profundamente transformadora:

\begin{quote}
  \textbf{Nunca escreva código de produção antes de ter um teste que
  falhe por sua ausência.}
\end{quote}

\indice{Test-Driven Development} \indice{Kent Beck}

O ciclo TDD opera em três estágios, ilustrados na Figura~\ref{fig:tdd-cycle}:

\begin{figure}[H]
  \centering
  \begin{tikzpicture}[
    node distance=2.2cm,
    redbox/.style={rectangle, draw=corTesteFalha, fill=corTesteFalha!10, text width=4cm, align=center, align=center, minimum height=1.8cm, rounded corners=4pt},
    greenbox/.style={rectangle, draw=corTesteSucesso, fill=corTesteSucesso!10, text width=4cm, align=center, align=center, minimum height=1.8cm, rounded corners=4pt},
    bluebox/.style={rectangle, draw=corSecundaria, fill=corSecundaria!10, text width=4cm, align=center, align=center, minimum height=1.8cm, rounded corners=4pt},
    arrow/.style={->, >=stealth, thick},
  ]
    \node[align=center, redbox] (red) {\textbf{[ICON] RED}\\\footnotesize Escrever teste que falha\\\footnotesize (contrato não atendido)};
    \node[align=center, greenbox, right=of red] (green) {\textbf{[ICON] GREEN}\\\footnotesize Código mínimo para\\\footnotesize o teste passar};
    \node[align=center, bluebox, below=of green] (refactor) {\textbf{[ICON] REFACTOR}\\\footnotesize Melhorar código\\\footnotesize mantendo testes verdes};

    \draw[arrow] (red) -- (green);
    \draw[arrow] (green) -- (refactor);
    \draw[arrow, dashed, corDestaque] (refactor) -| (red) node[near start, below, corDestaque] {\footnotesize próximo teste};
  \end{tikzpicture}
  \caption{Ciclo TDD: RED (falha) $\rightarrow$ GREEN (passa) $\rightarrow$ REFACTOR (melhora).}
  \label{fig:tdd-cycle}
\end{figure}

\subsection{RED: O Teste Que Falha}

O estágio RED é \textbf{contraintuitivo mas essencial}: você escreve um
teste que \textit{sabidamente} vai falhar, porque o código que ele testa
ainda não existe (ou ainda não atende ao critério). Este teste funciona
como \textbf{contrato}: ele define, de forma executável, o que o código
\textit{deve} fazer.

\begin{pratica}{RED: Teste antes da implementação}
  Escreva o teste que define o contrato da função \texttt{preprocess\_dental\_image}
  antes de implementá-la.
\end{pratica}

\begin{lstlisting}[language=Python, caption=Teste RED — define o contrato antes da implementação]
# TESTE (escrito PRIMEIRO — vai falhar!)
def test_preprocess_output_shape():
    """A saída deve ter dimensão (224, 224) após redimensionamento."""
    img = np.random.randint(0, 255, (300, 400), dtype=np.uint8)
    result = preprocess_dental_image(img)  # Função ainda não existe!
    assert result.shape == (224, 224)
# pytest: FAILED — NameError: name 'preprocess_dental_image' is not defined
\end{lstlisting}

\testefalha{O teste RED falha com \texttt{NameError}: a função ainda não
foi implementada. Este é o comportamento esperado e desejável no estágio RED.}

\subsection{GREEN: O Código Mínimo Que Passa}

No estágio GREEN, você escreve \textbf{o código mínimo necessário} para
fazer o teste passar — nada mais. Se o teste verifica apenas o \textit{shape},
basta uma função que redimensione a imagem. Não adicione CLAHE, normalização
ou qualquer outra feature que o teste não exija.

\begin{lstlisting}[language=Python, caption=GREEN — implementação mínima]
def preprocess_dental_image(image):
    """Redimensiona imagem odontológica para 224x224."""
    import cv2
    return cv2.resize(image, (224, 224))
# pytest: PASSED [ICON]
\end{lstlisting}

\teste{O teste GREEN passa: \texttt{result.shape == (224, 224)}. Implementação
mínima atende ao contrato.}

\subsection{REFACTOR: Melhorar Sem Quebrar}

Com o teste verde, você pode \textbf{refatorar} o código para melhorar
legibilidade, performance ou generalidade — desde que o teste continue
passando. Neste estágio, adicionamos CLAHE, normalização e validação.

\codigo{codigos/cap13/tdd_dental_testdata.py}

\teste{Após refatoração, todos os testes originais continuam passando.
O \texttt{preprocess\_batch} foi adicionado como otimização, sem quebrar
os testes existentes.}

\section{pytest e unittest no Google Colab}

O ecossistema Python oferece duas ferramentas principais para TDD:
\texttt{unittest} (biblioteca padrão) e \texttt{pytest} (terceiro, mais
expressivo). Para desenvolvimento odontológico, recomendamos \texttt{pytest}
por sua sintaxe concisa e rica coleção de fixtures e plugins.

\subsection{Configuração no Google Colab}

\begin{lstlisting}[language=bash, caption=Instalação e configuração do pytest no Colab]
!pip install pytest pytest-cov pytest-xdist
!mkdir -p /content/tests
\end{lstlisting}

\begin{lstlisting}[language=Python, caption=Estrutura básica de teste com pytest no Colab]
# %%writefile /content/tests/test_preprocessing.py
import numpy as np
import pytest

class TestDentalPreprocessing:
    """Suíte de testes para pré-processamento de imagens odontológicas."""

    def test_resize_preserves_aspect_ratio(self):
        """Redimensionamento deve preservar razão de aspecto."""
        img = np.random.randint(0, 255, (300, 400), dtype=np.uint8)
        # Implementação deve manter proporção
        assert img.shape[1] / img.shape[0] == 400/300

    def test_normalization_range(self):
        """Normalização deve mapear pixels para [0, 1]."""
        img = np.random.randint(0, 255, (224, 224), dtype=np.uint8)
        normalized = img.astype(np.float32) / 255.0
        assert normalized.min() >= 0.0
        assert normalized.max() <= 1.0

    @pytest.mark.parametrize("size", [(128, 128), (256, 256), (512, 512)])
    def test_multiple_resolutions(self, size):
        """Pré-processamento deve funcionar em múltiplas resoluções."""
        img = np.random.randint(0, 255, size, dtype=np.uint8)
        # Verifica que o processamento não lança exceção
        assert img.shape == size
\end{lstlisting}

\subsection{Executando Testes no Colab}

\begin{lstlisting}[language=bash, caption=Execução de testes com pytest]
!python -m pytest /content/tests/ -v --tb=short
# ============================= test session starts =============================
# collected 3 items
# tests/test_preprocessing.py::TestDentalPreprocessing::test_resize_preserves [ICON]
# tests/test_preprocessing.py::TestDentalPreprocessing::test_normalization [ICON]
# tests/test_preprocessing.py::TestDentalPreprocessing::test_multiple_resolutions[128] [ICON]
# tests/test_preprocessing.py::TestDentalPreprocessing::test_multiple_resolutions[256] [ICON]
# tests/test_preprocessing.py::TestDentalPreprocessing::test_multiple_resolutions[512] [ICON]
# ============================= 5 passed in 0.15s ==============================
\end{lstlisting}

\notaexplicativa{O Google Colab permite executar comandos de terminal
prefixados com \texttt{!}. Isso significa que você pode rodar \texttt{pytest}
diretamente de uma célula do notebook, integrando TDD ao fluxo de trabalho
interativo.}

\section{Fixtures para Dados Odontológicos Sintéticos}

Uma \textit{fixture} é um mecanismo do pytest que fornece dados de teste
consistentes e reutilizáveis. O pacote \texttt{odontoia.tdd.test\_helpers}
oferece fixtures especializadas para o domínio odontológico.

\subsection{DentalTestData: Estrutura de Dados de Teste}

A classe \texttt{DentalTestData} é o contêiner central para dados
odontológicos sintéticos:

\begin{description}
  \item[\texttt{images}] ndarray (N, H, W) — radiografias sintéticas.
  \item[\texttt{labels}] ndarray (N,) — rótulos de classe.
  \item[\texttt{masks}] ndarray (N, H, W) opcional — máscaras de segmentação.
  \item[\texttt{patient\_ids}] List[str] opcional — identificadores de pacientes.
  \item[\texttt{feature\_names}] List[str] opcional — nomes das colunas (tabular).
  \item[\texttt{feature\_matrix}] ndarray (N, F) opcional — dados tabulares.
\end{description}

Propriedades computadas:
\begin{itemize}
  \item \texttt{n\_samples} — número de amostras.
  \item \texttt{n\_classes} — número de classes únicas.
  \item \texttt{image\_shape} — tupla (H, W) das imagens.
  \item \texttt{split(test\_size, random\_state)} — divide em treino/teste.
\end{itemize}

\subsection{Geradores de Dados Sintéticos}

O módulo \texttt{test\_helpers} fornece três geradores principais:

\subsubsection{Radiografias sintéticas}

\textbf{\texttt{create\_test\_radiograph(size, n\_classes, random\_state)}}
— Gera um dataset de radiografias sintéticas com estruturas elípticas
simulando dentes. Ideal para testar pipelines de classificação de imagens.

\subsubsection{Dados de paciente odontológico}

\textbf{\texttt{create\_test\_patient(patient\_id, idade, fumante, diabetes, periodontal\_stage)}}
— Cria um registro de paciente odontológico com dados clínicos simulados:
32 dentes (notação FDI), profundidade de sondagem, mobilidade e exames.

\subsubsection{Dataset periodontal tabular}

\textbf{\texttt{generate\_periodontal\_dataset(n\_samples, random\_state)}}
— Gera dataset tabular com 10 features periodontais: profundidade de
sondagem, nível de inserção, sangramento, perda óssea, mobilidade, placa,
fumante, diabetes, idade e higiene oral.

\begin{pratica}{Criando datasets sintéticos para TDD odontológico}
  Objetivo: Utilizar os geradores do \texttt{odontoia.tdd.test\_helpers}
  para criar dados de teste e validar pipelines antes de usar dados reais.
\end{pratica}

\teste{\texttt{create\_test\_radiograph((256,256), n\_classes=3)} gera 32 imagens
de $256\times256$ com 3 classes, valores no intervalo [0, 255], e método
\texttt{split()} funcional.}

\teste{\texttt{create\_test\_patient("PAC-042", idade=45)} gera registro
completo com 32 dentes, exames periodontais e dados demográficos.}

\section{Testes para Pré-processamento de Imagens Odontológicas}

O pré-processamento é a etapa mais crítica — e mais propensa a erros
sutis — em qualquer pipeline de IA odontológica. Testes rigorosos nesta
etapa previnem que erros de normalização, redimensionamento ou equalização
se propaguem para o modelo, produzindo resultados clinicamente inválidos.

\subsection{Testes de Dimensionalidade}

\begin{lstlisting}[language=Python, caption=Testes de dimensionalidade para pré-processamento]
def test_resize_output_shape():
    """Redimensionamento deve produzir exatamente o tamanho solicitado."""
    from odontoia.tdd.test_helpers import create_test_radiograph
    data = create_test_radiograph(size=(300, 400), n_classes=2)
    img = data.images[0]
    import cv2
    resized = cv2.resize(img, (224, 224))
    assert resized.shape == (224, 224)

def test_grayscale_preserved():
    """Imagem grayscale deve permanecer 2D após pré-processamento."""
    img = np.random.randint(0, 255, (256, 256), dtype=np.uint8)
    processed = img.astype(np.float32) / 255.0
    assert processed.ndim == 2  # Não deve ganhar canal de cor
\end{lstlisting}

\subsection{Testes de Normalização}

\begin{lstlisting}[language=Python, caption=Testes de normalização de intensidade]
def test_normalization_bounds():
    """Normalização [0,1] deve preservar extremos relativos."""
    img = np.array([[0, 128, 255]], dtype=np.uint8)
    normalized = img.astype(np.float32) / 255.0
    assert normalized.min() == 0.0
    assert normalized.max() == 1.0
    assert normalized[0, 1] == pytest.approx(128/255, abs=1e-4)

def test_normalization_idempotent():
    """Normalizar duas vezes não deve alterar o resultado."""
    img = np.random.randint(0, 255, (100, 100), dtype=np.uint8)
    n1 = img.astype(np.float32) / 255.0
    n2 = (n1 * 255).astype(np.uint8).astype(np.float32) / 255.0
    # Pequena perda por conversão uint8 → float32
    assert np.allclose(n1, n2, atol=1/255)
\end{lstlisting}

\subsection{Testes de CLAHE e Equalização}

\teste{\texttt{createCLAHE(clipLimit=2.0, tileGridSize=(8,8))}
melhora o contraste local sem introduzir artefatos: PSNR $>$ 30 dB
entre original e equalizada.}

\testefalha{CLAHE com \texttt{clipLimit=0} ou \texttt{tileGridSize=(1,1)}
produz resultado igual à imagem original — não há equalização efetiva.}

\section{Testes para Métricas Clínicas}

As métricas clínicas são o critério último de qualidade de um sistema
de IA odontológica. O módulo \texttt{odontoia.tdd.test\_helpers} fornece
a função \texttt{assert\_clinical\_metrics} para validação automatizada.

\subsection{assert\_clinical\_metrics}

\begin{lstlisting}[language=Python, caption={Uso de assert\_clinical\_metrics em TDD}]
from odontoia.tdd.test_helpers import assert_clinical_metrics

def test_model_meets_clinical_thresholds():
    """Modelo deve atender aos limiares clínicos mínimos."""
    y_true = np.array([0, 0, 0, 0, 1, 1, 1, 1, 1, 1])
    y_pred = np.array([0, 0, 0, 0, 1, 1, 1, 1, 1, 1])  # Perfeito

    result = assert_clinical_metrics(
        y_true, y_pred,
        min_accuracy=0.85,
        min_sensitivity=0.80,
        min_specificity=0.80,
        raise_on_fail=False,
    )
    assert result["approved"] is True
    assert result["accuracy"] == 1.0
    assert result["sensitivity"] == 1.0
    assert result["specificity"] == 1.0
\end{lstlisting}

\teste{Com predições perfeitas, \texttt{assert\_clinical\_metrics} retorna
\texttt{approved=True} com todas as métricas em 1.0.}

\testefalha{Com predições aleatórias (50\% de acerto), \texttt{approved=False}
e a mensagem de erro lista quais métricas ficaram abaixo do limiar.}

\subsection{Métricas Suportadas}

A Tabela~\ref{tab:clinical-metrics-supported} lista as métricas suportadas
pelo \texttt{assert\_clinical\_metrics}:

\begin{table}[H]
  \centering
  \caption{Métricas clínicas verificáveis automaticamente.}
  \label{tab:clinical-metrics-supported}
  \begin{tabularx}{\textwidth}{lXl}
    \toprule
    \textbf{Métrica} & \textbf{Fórmula / Descrição} & \textbf{Limiar Padrão} \\
    \midrule
    Acurácia & $(TP + TN) / (TP + TN + FP + FN)$ & $>$ 0,80 \\
    Sensibilidade & $TP / (TP + FN)$ & $>$ 0,75 \\
    Especificidade & $TN / (TN + FP)$ & $>$ 0,75 \\
    Precisão & $TP / (TP + FP)$ & $>$ 0,75 \\
    F1-Score & $2 \times (P \times S) / (P + S)$ & $>$ 0,75 \\
    \bottomrule
  \end{tabularx}
\end{table}

\notaexplicativa{Para classificação multiclasse, as métricas são calculadas
como média macro (\textit{one-vs-rest}), garantindo que classes minoritárias
(como carcinoma espinocelular) não sejam ignoradas na avaliação.}

\section{Mocking de Datasets e Modelos}

Em TDD, \textit{mocking} é a técnica de substituir dependências reais
(como datasets de 50 GB ou modelos que levam horas para treinar) por
versões simplificadas que retornam respostas pré-definidas.

\subsection{Mock de Dataset DICOM}

\begin{lstlisting}[language=Python, caption=Mock para carregador de imagens DICOM]
from unittest.mock import patch, MagicMock
import numpy as np

def test_dicom_loader_with_mock():
    """Testa o pipeline sem acessar arquivos DICOM reais."""
    mock_dicom = np.random.randint(0, 4096, (512, 512), dtype=np.uint16)

    with patch("pydicom.dcmread") as mock_read:
        mock_dataset = MagicMock()
        mock_dataset.pixel_array = mock_dicom
        mock_read.return_value = mock_dataset

        # Agora o teste executa sem arquivos reais
        from odontoia.data.loader import load_dicom_as_array
        result = load_dicom_as_array("fake/path.dcm")
        assert result.shape == (512, 512)
        assert result.dtype == np.uint16
\end{lstlisting}

\subsection{Mock de Modelo Treinado}

\begin{lstlisting}[language=Python, caption=Mock de modelo para testar pipeline de inferência]
def test_inference_pipeline_with_mock_model():
    """Testa o pipeline de inferência sem carregar pesos reais."""
    mock_model = MagicMock()
    mock_model.predict.return_value = np.array([[0.1, 0.8, 0.1]])

    # Testa o pós-processamento da predição
    probs = mock_model.predict(np.zeros((1, 224, 224, 3)))
    predicted_class = np.argmax(probs, axis=1)[0]
    confidence = np.max(probs)

    assert predicted_class == 1
    assert confidence == 0.8
    assert probs.shape == (1, 3)
\end{lstlisting}

\teste{O mock do modelo permite testar a lógica de pós-processamento
(pipeline de inferência) sem depender de um modelo real treinado,
reduzindo o tempo de execução do teste de horas para milissegundos.}

\subsection{Mock de API Externa}

Para sistemas que consultam APIs externas (ex.: Hugging Face Inference
API, DICOMweb), o mocking é essencial para testes offline:

\begin{lstlisting}[language=Python, caption=Mock de API externa para teste offline]
from unittest.mock import patch
import json

def test_api_inference_with_mock():
    """Testa integração com API sem conexão de rede."""
    mock_response = MagicMock()
    mock_response.status_code = 200
    mock_response.json.return_value = {
        "predictions": [{"class": "carie", "confidence": 0.92}]
    }

    with patch("requests.post", return_value=mock_response):
        import requests
        resp = requests.post("https://api.hf.co/models/odontologia/caries",
                             json={"image": "base64..."})
        assert resp.status_code == 200
        data = resp.json()
        assert data["predictions"][0]["class"] == "carie"
        assert data["predictions"][0]["confidence"] > 0.9
\end{lstlisting}

\section{Cobertura de Testes e CI/CD}

\subsection{Cobertura de Código com pytest-cov}

A cobertura de código (\textit{code coverage}) mede qual porcentagem do
código-fonte é exercitada pelos testes. Embora 100\% de cobertura não
garanta ausência de bugs, \textbf{baixa cobertura é um forte indicador
de risco}.

\begin{lstlisting}[language=bash, caption=Relatório de cobertura com pytest-cov]
!python -m pytest tests/ -v \
    --cov=odontoia \
    --cov-report=term-missing \
    --cov-report=html \
    --cov-fail-under=80
\end{lstlisting}

\destaque{Para software odontológico com potencial uso clínico-assistido,
recomendamos cobertura mínima de 85\%. Módulos de inferência e
pré-processamento devem ter cobertura $\geq$ 95\%.}

\subsection{CI/CD com GitHub Actions}

A integração contínua (CI) automatiza a execução de testes a cada
\texttt{git push}, garantindo que nenhuma regressão passe despercebida.

\begin{lstlisting}[language=yaml, caption=GitHub Actions workflow para testes odontológicos]
# .github/workflows/tests.yml
name: Dental AI Tests
on: [push, pull_request]
jobs:
  test:
    runs-on: ubuntu-latest
    steps:
      - uses: actions/checkout@v4
      - uses: actions/setup-python@v5
        with:
          python-version: '3.10'
      - name: Install dependencies
        run: |
          pip install -e .
          pip install pytest pytest-cov
      - name: Run TDD tests
        run: |
          python -m pytest tests/ -v --cov=odontoia \
            --cov-fail-under=80 \
            --junitxml=test-results.xml
      - name: Verify SDD specs
        run: |
          python -c "from odontoia.sdd.validator import verify_all_specs; \
                     results = verify_all_specs(); \
                     assert all(r.valid for r in results.values())"
\end{lstlisting}

\teste{O workflow CI/CD acima executa \textbf{duas camadas de verificação}:
(1) testes TDD via pytest com cobertura mínima de 80\% e (2) validação
de todas as especificações SDD registradas via SpecVerifier.}

\subsection{Pirâmide de Testes Odontológicos}

A Figura~\ref{fig:test-pyramid} adapta a clássica pirâmide de testes
para o contexto odontológico:

\begin{figure}[H]
  \centering
  \begin{tikzpicture}[scale=0.9]
    % Camada superior (UI / E2E)
    \fill[corDestaque!20] (0,0) -- (8,0) -- (7,3) -- (1,3) -- cycle;
    \draw[corDestaque, thick] (0,0) -- (8,0) -- (7,3) -- (1,3) -- cycle;
    \node at (4,1.2) {\begin{tabular}{c}\textbf{E2E Clínico}\\\footnotesize Workflow completo\\\footnotesize (5\% dos testes)\end{tabular}};

    % Camada intermediária (Integração)
    \fill[corSecundaria!15] (0.8,0) -- (7.2,0) -- (6.5,5.5) -- (1.5,5.5) -- cycle;
    \draw[corSecundaria, thick] (0.8,0) -- (7.2,0) -- (6.5,5.5) -- (1.5,5.5) -- cycle;
    \node at (4,3.5) {\begin{tabular}{c}\textbf{Integração}\\\footnotesize Pipeline pré-processamento\\\footnotesize + modelo + pós-processamento\\\footnotesize (15\% dos testes)\end{tabular}};

    % Camada base (Unidade)
    \fill[corTesteSucesso!12] (1.5,0) -- (6.5,0) -- (6,8) -- (2,8) -- cycle;
    \draw[corTesteSucesso, thick] (1.5,0) -- (6.5,0) -- (6,8) -- (2,8) -- cycle;
    \node at (4,5.5) {\begin{tabular}{c}\textbf{Unidade}\\\footnotesize Funções isoladas,\\fixtures, métricas clínicas\\\footnotesize (80\% dos testes)\end{tabular}};

    % Rótulo
    \node[right] at (8.5, 4) {\footnotesize Lento, frágil};
    \node[right, corDestaque] at (8.5, 1.2) {\footnotesize $\uparrow$};
    \node[right, corTesteSucesso] at (8.5, 5.5) {\footnotesize $\downarrow$};
    \node[right] at (8.5, 7) {\footnotesize Rápido, robusto};
  \end{tikzpicture}
  \caption{Pirâmide de testes odontológicos: 80\% unitários, 15\% integração, 5\% E2E.}
  \label{fig:test-pyramid}
\end{figure}

\section{Síntese do Capítulo}

Neste capítulo, cobrimos a aplicação completa do TDD ao desenvolvimento
de software odontológico. Os principais pontos são:

\begin{itemize}
  \item \textbf{Ciclo RED $\rightarrow$ GREEN $\rightarrow$ REFACTOR}:
    teste falha $\rightarrow$ código mínimo $\rightarrow$ melhoria segura.

  \item \textbf{pytest no Google Colab}: testes automatizados executam
    diretamente no notebook, sem sair do ambiente de desenvolvimento.

  \item \textbf{DentalTestData e fixtures sintéticas}: permitem testar
    pipelines odontológicos sem depender de datasets reais (grandes,
    privados, regulados).

  \item \textbf{assert\_clinical\_metrics}: valida automaticamente se
    métricas clínicas (acurácia, sensibilidade, especificidade) atendem
    aos limiares clinicamente significativos.

  \item \textbf{Mocking}: isola dependências externas (DICOM, APIs, modelos
    grandes), tornando os testes rápidos e determinísticos.

  \item \textbf{CI/CD + cobertura}: GitHub Actions executa testes e
    valida especificações a cada push, com cobertura mínima de 80\%.
\end{itemize}

O SDD (Capítulo~12) e o TDD (este capítulo) formam a dupla metodológica
que guiará todos os projetos práticos do Capítulo~14. Antes de implementar,
especifique; antes de considerar pronto, teste.
